\documentclass[11pt,a4paper]{article}
\usepackage[utf8]{inputenc}
\usepackage[T1]{fontenc}
\usepackage[polish]{babel}
\usepackage{lmodern}
\usepackage[margin=2.3cm]{geometry}
\usepackage{booktabs}
\usepackage{array}
\usepackage{enumitem}
\usepackage{titlesec}
\usepackage{xcolor}
\usepackage{fancyhdr}
\usepackage{hyperref}

\definecolor{navy}{RGB}{27,42,86}
\definecolor{lgray}{RGB}{240,240,244}

\titleformat{\section}{\large\bfseries\color{navy}}{\thesection.}{0.5em}{}
\titleformat{\subsection}{\bfseries\color{navy}}{\thesubsection.}{0.5em}{}

\pagestyle{fancy}
\fancyhf{}
\renewcommand{\headrulewidth}{0.6pt}
\fancyhead[L]{\small Neurofizjologia i pomiar EEG}
\fancyhead[R]{\small Instrukcja laboratoryjna nr 3}
\fancyfoot[C]{\small\thepage}

\newcommand{\infobox}[1]{%
  \par\noindent\colorbox{lgray}{\parbox{\dimexpr\linewidth-2\fboxsep}{\vspace{2pt}#1\vspace{2pt}}}\par\vspace{4pt}}

\title{\vspace{-1.5cm}\textbf{\color{navy}Instrukcja laboratoryjna nr 3}\\[4pt]
\Large Projektowanie eksperymentu EEG i rejestracja danych}
\author{Neurofizjologia i pomiar EEG\\ Psychologia, studia II stopnia --- specjalność neuropsychologia kliniczna}

\begin{document}
\maketitle
\thispagestyle{fancy}

\section{Cele zajęć}
\begin{itemize}[nosep]
  \item Nauka projektowania paradygmatu eksperymentalnego adekwatnego do pytania badawczego EEG.
  \item Zapoznanie z zasadami doboru i synchronizacji bodźców (znaczniki/triggery) z zapisem EEG.
  \item Ćwiczenie praktycznego przeprowadzenia rejestracji EEG w warunkach laboratoryjnych.
  \item Utrwalenie zasad prawidłowego kontaktu z osobą badaną w trakcie właściwego badania.
\end{itemize}

\section{Wprowadzenie teoretyczne}

\subsection{Rodzaje paradygmatów badawczych w EEG}
\begin{itemize}[nosep]
  \item \textbf{Rejestracja spoczynkowa (resting-state)} -- z oczami otwartymi (EO) i zamkniętymi (EC), 3--5 minut na warunek; baza do analiz mocy widmowej.
  \item \textbf{Paradygmaty związane ze zdarzeniami (event-related)} -- prezentacja bodźców z rejestrowanymi znacznikami czasowymi, podstawa analizy potencjałów wywołanych (ERP).
  \item \textbf{Paradygmaty typu oddball} -- rzadkie bodźce docelowe wśród częstych standardowych, stosowane do badania komponentu P300.
  \item \textbf{Zadania ciągłe} -- np. zadania Go/No-Go, Stroop -- łączą pomiar behawioralny (czas reakcji, poprawność) z zapisem EEG.
\end{itemize}

\subsection{Synchronizacja bodźców i znaczników (triggerów)}
Precyzyjna analiza ERP wymaga dokładnego, submilisekundowego znakowania momentu prezentacji bodźca w strumieniu EEG. Realizuje się to poprzez wysyłanie kodów cyfrowych (\emph{event markers}) z komputera prezentującego bodźce do wzmacniacza EEG, zwykle przez port równoległy/USB lub interfejs sieciowy dedykowanego oprogramowania. Błędy synchronizacji (jitter) obniżają precyzję pomiaru latencji komponentów ERP.

\subsection{Kontrola zmiennych zakłócających}
Projekt eksperymentu EEG musi minimalizować artefakty przewidywalne już na etapie planowania: ruchy gałek ocznych i mrugnięcia (instrukcja dot. fiksacji wzroku), napięcie mięśniowe (wygodna pozycja, oparcie głowy), zakłócenia elektryczne otoczenia (ekranowanie, uziemienie sprzętu).

\infobox{\textbf{Zasada etyczna:} Osoba badana musi wyrazić świadomą zgodę, być poinformowana o przebiegu badania i mieć prawo do rezygnacji na każdym etapie bez podawania przyczyny.}

\section{Sprzęt i oprogramowanie}
Stanowisko EEG z poprzednich zajęć, komputer do prezentacji bodźców (np. PsychoPy/E-Prime), kabel/interfejs do wysyłania znaczników, gotowy skrypt prostego paradygmatu oddball lub zadania Go/No-Go.

\section{Przebieg zajęć (harmonogram, 120 min)}
\begin{center}
\begin{tabular}{@{}p{2.3cm}p{11.7cm}@{}}
\toprule
\textbf{Czas} & \textbf{Aktywność}\\
\midrule
0--15 min & Wykład: rodzaje paradygmatów, znaczniki czasowe, kontrola artefaktów.\\
15--35 min & Prezentacja przykładowego skryptu paradygmatu oddball i konfiguracji znaczników.\\
35--55 min & Praca w grupach: projektowanie własnego prostego paradygmatu (zadanie 1).\\
55--65 min & Przerwa.\\
65--90 min & Przygotowanie osoby badanej (czepek, impedancja) i próbna rejestracja EEG (zadanie 2).\\
90--115 min & Właściwa rejestracja krótkiego paradygmatu w rolach naprzemiennych badacz/osoba badana (zadanie 3).\\
115--120 min & Zapis i zabezpieczenie danych, podsumowanie.\\
\bottomrule
\end{tabular}
\end{center}

\section{Zadania do wykonania}
\begin{enumerate}[label=\textbf{Zadanie \arabic*.}, leftmargin=2.7cm]
  \item W grupach zaprojektujcie krótki paradygmat oddball (bodźce standardowe 80\%, docelowe 20\%) wraz z instrukcją dla osoby badanej i schematem kodowania znaczników.
  \item Przygotujcie stanowisko badawcze zgodnie z procedurą z zajęć 2 i wykonajcie 2-minutową rejestrację próbną w celu weryfikacji jakości sygnału i poprawności znaczników.
  \item Przeprowadźcie właściwą, ok. 8-minutową rejestrację zaprojektowanego paradygmatu, zamieniając się rolami badacza i osoby badanej.
\end{enumerate}

\section{Kryteria zaliczenia ćwiczenia}
Poprawność projektu paradygmatu (30\%), jakość techniczna przeprowadzonej rejestracji, w tym zsynchronizowanie znaczników (40\%), właściwy kontakt z osobą badaną (30\%).

\section{Pytania kontrolne}
\begin{enumerate}[nosep]
  \item Czym różni się paradygmat spoczynkowy od paradygmatu związanego ze zdarzeniami?
  \item Dlaczego precyzja znaczników czasowych jest krytyczna dla analizy ERP?
  \item Jakie źródła artefaktów można ograniczyć już na etapie projektowania badania?
  \item Jakie elementy powinna zawierać instrukcja dla osoby badanej w zadaniu oddball?
\end{enumerate}

\section{Literatura}
\begin{enumerate}[nosep]
  \item Szumska, I. (2019). \emph{Neurofizjologia i neuropsychologia -- skrypt dla studentów AEH}. Wydawnictwo AEH.
  \item Chaddad, A., Wu, Y., Kateb, R., \& Bouridane, A. (2023). Electroencephalography signal processing. \emph{Sensors, 23}(14), 6434.
  \item Kalat, J. W. (2019). \emph{Biological Psychology}. Cengage Learning.
\end{enumerate}

\end{document}
